\documentclass{beamer}
\usepackage{ctex}
\usepackage{tikz}
\usepackage{hologo}
\usepackage{cncolours}
\usepackage{amsthm,amsmath}
\usepackage{fancyvrb}
\usepackage{url}

\usetikzlibrary{calc}
\usetheme{WHUSpot}
\usefonttheme{serif}
\mode<presentation>

\title{武汉大学主题 Beamer 模板 - WHUSpot}
\author{向飞宇}
\institute{武汉大学}
\date{\today}

\begin{document}

\begin{frame}
    \frametitle{\\}
    \titlepage
\end{frame}

\section{使用方法}

\begin{frame}
    \frametitle{使用方法与注意事项}
    \begin{itemize}
        \item 用 $\backslash \mbox{usetheme\{WHUSpot\}}$ 即可载入模板，当然也需要将 beamerthemeWHUSpot.sty 放在合适的地方（其实在同目录下就行）.
        \item 请使用 \hologo{XeLaTeX} 编译，一般而言使用\alert{TeX Live 2018}及以上版本能够保证正常编译（当然稍微低版本的TeX Live应该也没问题）.
        \item 本人\alert{没有设计subsection及以下层次}的页面，仅在有section时调用一个特别的页面，请谨慎组织标题内容.
        \item 请不要轻易改变配色方案.
    \end{itemize}
\end{frame}

\begin{frame}
    \frametitle{列表与枚举环境}
    \begin{itemize}
        \item 列表环境是beamer默认的，可以修改. \pause
        \item 支持多级列表，但修改环境要逐级修改，\alert{不建议}使用. \pause
    \end{itemize}

    \begin{enumerate}
        \item 枚举环境是beamer默认的，可以修改. \pause
        \item 支持多级枚举，但修改环境要逐级修改，\alert{不建议}使用.
    \end{enumerate}
\end{frame}

\begin{frame}
    \frametitle{定理的一个例子}
    \begin{theorem}[Risez定理的一个应用]
    \begin{equation*}
        f_n \stackrel{m}{\longrightarrow} f \Leftrightarrow \forall \{f_{n_k}\} \subset \{f_n\}, \exists \{f_{n_{k'}}\} \subset \{f_{n_k}\}, s.t. \quad f_{n_{k'}} \to f \quad a.e.     
    \end{equation*}
    \end{theorem}
    \renewcommand{\qedsymbol}{}
    \begin{proof}
        （$\Rightarrow$）设$f_n \stackrel{m}{\longrightarrow}f$，则$\forall \{f_{n_k}\} \subset \{f_n\}$，$f_{n_k} \stackrel{m}{\longrightarrow} f$（直接按照定义证明），由\textbf{定理$3$}（Risez定理），$\exists \{f_{n_{k'}}\} \subset \{f_{n_k}\}, s.t. \quad f_{n_{k'}} \to f \quad a.e. (k' \to \infty).$ \\
        （$\Leftarrow$）用反证法，若$\{f_n\} \stackrel{m}{\nrightarrow} f$，$\exists \varepsilon > 0, s.t. \quad m(E(|f_n - f| \geqslant \varepsilon)) \nrightarrow 0$，于是$\exists \delta > 0, \{f_{n_k}\} \subset \{f_n\}, s.t.$ \[m(E(|f_{n_k} - f| \geqslant \varepsilon)) \geqslant \delta, \quad k = 1, 2, \cdots.\] 
    \end{proof}
\end{frame}

\begin{frame}
    \frametitle{定理的一个例子(续)}
    \begin{proof}[证明（续）]
        另一方面，由假设条件，$\exists \{f_{n_{k'}}\} \subset \{f_{n_k}\}, s.t. \quad f_{n_{k'}} \to f \quad a.e.$ \\
        因为$m(E) < \infty$，由\textbf{定理$2$}，$f_{n_{k'}} \stackrel{m}{\longrightarrow} f$，从而得出矛盾.从而$f_n \stackrel{m}{\longrightarrow} f$.
    \end{proof}
\end{frame}

\section{设计灵感}
\begin{frame}
    \frametitle{配色}
    红，蓝，黑为主，绿，很浅的白为辅.

    用cncolours包配色.
\end{frame}

\begin{frame}
    \frametitle{一瓣瓣的樱花}
    \begin{figure}
    \centering
    \begin{tikzpicture}[scale=.8]
        \foreach \x in {1,...,5}            
        {
            \pause
            \filldraw[out=\sakuraoutangle,in=\sakurainangle,rotate=-(\x - 1) * 72,\sakuraoutercolor] (0,0) to (\sakuraxoffset,\sakurayoffset);
            \filldraw[out=180-\sakuraoutangle,in=-180-\sakurainangle,rotate=-(\x - 1) * 72,\sakuraoutercolor] (0,0) to (-\sakuraxoffset,\sakurayoffset);
            \filldraw[rotate=-(\x - 1) * 72,\sakuraoutercolor] (0,0) -- (-\sakuraxoffset,\sakurayoffset) -- (0,\sakurayoffset-\sakuraxoffset) -- (\sakuraxoffset,\sakurayoffset) -- (0,0);
            \fill[rotate=-(\x - 1) * 72-6,\sakurainnercolor,opacity=.65] (0,\sakurayoffset *0.25) circle (1.2pt);
            \fill[rotate=-(\x - 1) * 72+6,\sakurainnercolor,opacity=.65] (0,\sakurayoffset *0.25) circle (1.2pt);
            \fill[\sakurakernelcolor,opacity=.75] (0,0) circle (2.2pt);
        }
    \end{tikzpicture}
    \end{figure}
\end{frame}

\begin{frame}
    \frametitle{老图书馆不是一日建成的}
    \begin{figure}
    \centering
    \renewcommand{\buildingcolor}{blue}
    \begin{tikzpicture}[scale=.4]
        \useasboundingbox (-0.2, 0.5) rectangle (0.05, 10);
        \fill[\buildingcolor] (-1, 0) rectangle (1, 2);
        \draw[\buildingcolor] (-1.2, 0) rectangle (1.2, 2.2);
        \draw[line width=2.5pt,\buildingcolor] (-6, 0) -- (6, 0);
        \foreach \x in {1,...,5}
        {
            \draw[line width=1.5pt, \buildingcolor] (-4.5 - \x * 0.2, -\x / 3) -- (4.5 + \x * 0.2, -\x / 3);
        }
        \fill[\buildingcolor] (-12, -2) -- (12, -2) -- (11.5, -2.2) -- (-11.5, -2.2) -- cycle;
        \fill[\buildingcolor] (-11.3, -2.28) -- (11.3, -2.28) -- (10.8, -2.48) -- (-10.8, -2.48) -- cycle;
        \fill[\buildingcolor] (-6, -2.56) -- (6, -2.56) -- (5.5, -2.76) -- (-5.5, -2.76) -- cycle;
        \pause

        \foreach \i in {-1, 1} 
            \filldraw[\buildingcolor] (-5.8 * \i, 0) -- (-5.8 * \i, -0.3) -- (-5.5 * \i, -0.3) -- (-6.4 * \i, -1.8) -- (-11.4 * \i, -1.8) -- (-10.68 * \i, -0.6) -- (-10.5 * \i, -0.6) -- (-10.5 * \i, 3.5) -- (-6 * \i, 3.5) -- (-6 * \i, 0) -- cycle;
        \foreach \j in {8.97, 7.72, 6.47}
            \filldraw[rounded corners, white] (\j, 2.4) rectangle (\j+ 1, 3);
        \foreach \j in {-8.97, -7.72, -6.47}
            \filldraw[rounded corners, white] (\j, 2.4) rectangle (\j- 1, 3);
        \pause

        \filldraw[\buildingcolor] (-10.5, 3.9) rectangle (10.5, 4.6);
        \filldraw[\buildingcolor] (-10.5, 3.5) rectangle (10.5, 3.7);
        \foreach \k in {-5.6, -3.1, 3.5}
            \filldraw[\buildingcolor] (\k, 0) -- (\k, 3.9) -- (\k + 0.5, 3.9) -- (\k + 0.5, 3.5) -- (\k + 0.2, 3.3) -- (\k + 0.2, 0) -- cycle;
        \foreach \k in {5.6, 3.1, -3.5}
            \filldraw[\buildingcolor] (\k, 0) -- (\k, 3.9) -- (\k - 0.5, 3.9) -- (\k - 0.5, 3.5) -- (\k - 0.2, 3.3) -- (\k - 0.2, 0) -- cycle;
        \filldraw[\buildingcolor] (-0.25, 3.7) rectangle (0.25, 4);
        \pause

        \filldraw[\buildingcolor] (-10.5, 4.8) [out=165,in=-78]to (-12.8, 7) [out=-75,in=-180]to (-10.8, 6) [out=0,in=-95]to (-8.9, 8) [out=-85,in=-180]to (-6.8, 5.9) [out=0,in=-120]to (-5, 6.8) [out=-95,in=15]to (-6, 5.0) -- (-6, 4.8) -- cycle; 
        \filldraw[\buildingcolor] (10.5, 4.8) [out=15,in=-102]to (12.8, 7) [out=-105,in=0]to (10.8, 6) [out=180,in=-85]to (8.9, 8) [out=-95,in=0]to (6.8, 5.9) [out=180,in=-60]to (5, 6.8) [out=-85,in=165]to (6, 5.0) -- (6, 4.8) -- cycle;
        \filldraw[\buildingcolor] (-6, 4.8) rectangle (6, 5.0);
        \foreach \i in {-1, 1}
        {
            \filldraw[\buildingcolor] (-8.4 * \i, 6.9) -- (-5.1 * \i, 6.9) -- (-5.3 * \i, 6.6) -- (-8.1 * \i, 6.6) -- cycle;
            \filldraw[\buildingcolor] (-7.99 * \i, 6.45) -- (-5.4 * \i, 6.45) -- (-5.52 * \i, 6.35) -- (-7.87 * \i, 6.35) -- cycle;     
        }
        \filldraw[\buildingcolor] (-4.88, 6.9) -- (4.88, 6.9) -- (4.9, 6.6) -- (-4.9, 6.6) -- cycle;
        \filldraw[\buildingcolor] (-4.93, 6.45) -- (4.93, 6.45) -- (4.95, 6.35) -- (-4.95, 6.35) -- cycle;
        \pause

        \foreach \i in {-1, 1}
        {
            \filldraw[\buildingcolor] (-5.5 * \i, 5.1) -- (-5.5 * \i, 5.2) -- (-5.3 * \i, 5.3) -- (-5.3 * \i, 5.1) -- cycle;
            \filldraw[\buildingcolor] (-5.2 * \i, 5.1) -- (-5.2 * \i, 5.35) -- (-5.05 * \i, 5.85) -- (-5.05 * \i, 5.1) -- cycle;
        }
        \foreach \i in {0,...,24}
            \filldraw[\buildingcolor] (-4.9 + \i * 0.4, 5.1) rectangle (-4.9 + \i * 0.4 + 0.2, 6.25);
        \pause

        \filldraw[\buildingcolor] (-5, 7.7) [out=180,in=-90]to (-6.2, 8.5) [out=-30,in=-150]to (-4.4, 8.5) -- (-4.5, 9.6) to (-3.5, 8.8) -- (-2.7, 9.1) -- (-2.7, 10.2) -- (2.7, 10.2) -- (2.7, 9.1) -- (3.5, 8.8) to (4.5, 9.6) -- (4.4, 8.5) [out=-30,in=-150]to (6.2, 8.5) [out=-90,in=0]to (5, 7.7) -- (3.1, 8.2) -- (-3.1, 8.2) -- cycle;
        \filldraw[rounded corners=1pt,\buildingcolor] (-2.9, 10.0) rectangle (-2.5, 10.4);
        \filldraw[rounded corners=1pt,\buildingcolor] (2.9, 10.0) rectangle (2.5, 10.4);
        \foreach \i in {-1, 1}
            \filldraw[rounded corners=3pt,rotate around={13 * \i:(-4.7 * \i,7.1)},\buildingcolor] (-4.7 * \i, 7.1) rectangle (-3.1 * \i, 7.5); 
        \foreach \i in {0,...,3}
            \filldraw[rounded corners=3pt, \buildingcolor] (-2.95 + \i * 1.5, 7.5) rectangle (-1.55 + \i * 1.5, 7.9);
        \pause

        \filldraw[\buildingcolor] (-0.3, 10.2) rectangle (0.3, 10.3);
        \filldraw[\buildingcolor] (-0.3, 10.4) [bend left=10]to (-0.3, 11.2) -- (-0.2, 11.3) -- (-0.3, 11.4) -- (-0.3, 11.6) -- (0.3, 11.6) -- (0.3, 11.4) -- (0.2, 11.3) -- (0.3, 11.2) [bend left=10]to (0.3, 10.4) -- cycle;
        \filldraw[\buildingcolor] (-0.2, 11.7) rectangle (0.2, 11.8);
    \end{tikzpicture}
    \end{figure}
\end{frame}

\begin{frame}
    \frametitle{杂项}
    \begin{itemize}
        \item 本来想用樱花实现一个更复杂的progress bar的，最后放弃了.
        \item 因为背景用的老图书馆的背景，使用一个带颜色的block显得多余，所以很多地方直接使用了 beamer 的默认配置.
    \end{itemize}
\end{frame}

\begin{frame}
    \frametitle{鸣谢}
    \begin{itemize}
        \item 武汉大学.
        \item 用cncolours当调色板. \\ \url{https://github.com/liantze/pgfornament-han}
        \item progress bar的思路. \\ \url{https://github.com/DjalelBBZ/Algiers-beamer-template} 
    \end{itemize}
\end{frame}

\section{联系，bug反馈，授权}
\begin{frame}
    \frametitle{联系，bug反馈，授权}
    \begin{itemize}
        \item 最新版本在 \\ \url{https://github.com/T0nyX1ang/WHU-BeamerTemplate} 上发布.
        \item bug反馈可以直接在Github上发issue.
        \item 使用LPPL-1.3c（及以后版本）授权.
    \end{itemize}
\end{frame}

\end{document}